\documentclass[11pt]{article}

\usepackage[margin=1in]{geometry}
\usepackage{amsmath,amssymb}
\usepackage{booktabs}
\usepackage{array}
\usepackage{enumitem}
\usepackage[round,authoryear]{natbib}
\usepackage[hidelinks]{hyperref}
\usepackage{titlesec}
\usepackage{xcolor}

\titleformat{\section}{\large\bfseries}{\thesection.}{0.6em}{}
\titleformat{\subsection}{\normalsize\bfseries}{\thesubsection}{0.6em}{}
\setlist{itemsep=3pt,topsep=3pt}

\newcommand{\Mff}{M_{\mathrm{ff}}}
\newcommand{\IC}{\mathrm{IC}}
\newcommand{\Kint}{K^{\text{int}}}

\title{\bfseries A Beta-Neutral Long/Short Equity Strategy \\ in Software \& Services}
\author{Systematic Equities Research --- Project Plan}
\date{June 2026}

\begin{document}
\maketitle

%==================================================================
\section{Problem Statement}
%==================================================================

The objective is to design a systematic long/short equity portfolio, restricted to the
GICS Software \& Services industry group, that is neutral to directional equity beta and
that maximizes its out-of-sample, post-cost Sharpe ratio over a twenty-year forward
horizon. Directional market exposure is assumed hedged, so the portfolio is judged on the
quality of its relative bets; it remains free to take a net long or short tilt to the
industry where doing so improves the objective.

The universe is defined by the Software \& Services classification under GICS 2023.
Profiling the classification data gives roughly 1,100 securities over the full history,
split between Software and IT Services, dominated by the United States but with a
meaningful tail across Japan, the United Kingdom, Germany and a dozen other markets. The
point-in-time tradable set --- names that are alive, listed, and liquid enough to short
--- is far smaller, on the order of a few hundred stocks at any date.

Two features of this setting drive the design.

The first is limited breadth. By the fundamental law of active management,
$\mathrm{IR}\approx\IC\sqrt{\mathrm{BR}}$ \citep{grinold1989}, and a single narrow
industry group caps $\mathrm{BR}$. Performance therefore cannot rest on one factor; it
must come from combining weakly-correlated signals and raising $\IC$ with signals
tailored to software firms.

The second is that trading costs are steep and size-sensitive. The one-way cost in basis
points is
\begin{equation}
  c \;=\; 3\left(\frac{11}{\log_{10}\Mff}\right)^{6} + 3,
  \qquad
  \Mff = \max\!\big(\text{market cap} \times \text{FX} \times \text{free float},\;
  5\times 10^{7}\big),
\end{equation}
with $\Mff$ the freely-tradable market value in USD. The sixth power makes cost climb
sharply as size falls: roughly 6 bps at \$100B, 8 at \$10B, 13 at \$1B, 23 at \$100M, and
29 at the \$50M floor, paid twice per round-trip. After costs the strategy must favour
larger, liquid names and control turnover, and embedding that discipline in construction
likely contributes more to post-cost Sharpe than any additional signal.

The target is not a single high backtest number but a portfolio whose post-cost Sharpe
holds up under walk-forward testing and is stable across regions, sub-industries, and
regimes, built from a small set of orthogonal signals rather than a large, easily
over-fit collection.

%==================================================================
\section{Existing Research \& Approach}
%==================================================================

\subsection{Market factors}

The following are well-documented market factors and serve as benchmarks and building blocks.

\begin{itemize}
  \item \textbf{Value} \citep{fama1992}: earnings yield $E/P=\text{earnings}/\text{market
  cap}$; cheap stocks (high $E/P$) outperform. Long high, short low.

  \item \textbf{Momentum} \citep{jegadeesh1993}: cumulative total return over the prior
  twelve months, $\prod_{k=1}^{12}(1+r_{i,t-k})-1$. Long winners, short losers.

  \item \textbf{Short-term reversal} \citep{jegadeesh1990}: the most recent one-month total
  return $r_{i,t}$; recent losers outperform, so the book is long low, short high.

  \item \textbf{Profitability} \citep{novymarx2013}: gross profitability
  $(\text{revenue}-\text{COGS})/\text{assets}=\text{gross income}/\text{assets}$. Long
  profitable, short unprofitable.

  \item \textbf{Asset growth} \citep{cooper2008}: year-on-year asset growth
  $\text{assets}_t/\text{assets}_{t-12\text{m}}-1$; aggressive expanders underperform. Short
  high growth.

  \item \textbf{Net share issuance} \citep{pontiff2008}: year-on-year growth in diluted
  share count $\text{shares}_t/\text{shares}_{t-12\text{m}}-1$; issuers underperform,
  shrinkers outperform. Short positive issuance.

  \item \textbf{Low risk} \citep{frazzini2014}: trailing market beta
  $\beta_i=\operatorname{cov}(r_i,r_m)/\operatorname{var}(r_m)$; low-beta stocks earn higher
  risk-adjusted returns. Long low beta, short high; historically pronounced in technology.

  \item \textbf{Earnings momentum} \citep{bernard1989}: standardized unexpected earnings,
  the year-on-year change in trailing earnings scaled by its own trailing volatility,
  $\text{SUE}_{i,t}=(\text{E}_t-\text{E}_{t-12\text{m}})/\sigma_t(\Delta\text{E})$. Long
  positive surprises.

  \item \textbf{Earnings quality} \citep{sloan1996}: accruals
  $(\text{earnings}-\text{operating cash flow})/\text{assets}$; cash-backed earnings are
  more durable. Short high accruals.
\end{itemize}

\subsection{Software-tailored factors}

The following factors exploit how software firms report their economics, where standard accounting is most
misleading.

\begin{itemize}
  \item \textbf{Intangible-adjusted value and quality} \citep{petersTaylor2017}.
  Accounting expenses R\&D immediately rather than capitalizing it, so the reported book
  value and asset base of software firms are understated and ordinary value ratios
  misfire. We accumulate an intangible capital stock $\Kint$ from past R\&D by perpetual
  inventory,
  \begin{equation}
    \Kint_t \;=\; (1-\delta)\,\Kint_{t-1} \;+\; \text{R\&D}_t,
  \end{equation}
  where $\delta$ is the annual depreciation rate of the knowledge stock (taken near
  0.15--0.30). Adding $\Kint$ back into book equity and assets yields adjusted value and
  profitability ratios that compare software firms on consistent terms. R\&D productivity,
  $\Delta\text{sales}/\Kint$, captures whether research converts into growth.

  \item \textbf{Buyback quality} \citep{pontiff2008}. Repurchases that genuinely shrink the
  share count predict higher returns, but in software, where stock-based pay is the dominant
  source of dilution, reported buybacks often merely offset employee grants. The buyback
  reconciliation signal --- reported buyback yield minus the realized reduction in share
  count --- is negative for exactly those firms, flagging low-quality buybacks.
\end{itemize}

%==================================================================
\section{Proposed Experiments}
%==================================================================

\subsection{Experiment 1: Market factors across the three industries}

We evaluate the established market factors (Section 2.1) on each of the three permitted
industry universes --- Software \& Services, Banks \& Insurance, and Commodity Producers
--- by two complementary methods, then re-test under trading cost.

\textbf{Quintile sorts.} Each month, standardize every factor cross-sectionally relative to
its industry mean by z-score, rank the stocks, and split them into five equal quintiles.
Compare the average realized next-month return across the five quintiles; a monotonic
pattern indicates a genuine factor effect, and its slope and sign give the direction and
strength. Running this on all three universes shows whether an effect present in software
also holds in banks/insurance and commodity producers, or is industry-specific --- a guide
to how robust the software result is likely to be out-of-sample.

\textbf{Cross-sectional regressions.} Each month, regress the next-period cross-section of
stock returns on the factor \citep{famamacbeth1973}. Collect the resulting time series of
slope coefficients and report its t-statistic, computed with autocorrelation-consistent
standard errors, measuring whether the factor is significant across
time.

\textbf{Long/short Sharpe, gross and net of cost.} For each factor, form the portfolio that
is long the top quintile and short the bottom quintile (oriented by the factor's sign),
equal-weighted within each leg and rebalanced monthly, and report its gross Sharpe ratio.
Every change in position incurs a fractional cost, and the net portfolio return is therefore
\begin{equation}
  \tilde{r}^{\text{port}}_t \;=\; r^{\text{port}}_t \;-\;
  \sum_i \kappa_{i,t}\,\lvert w_{i,t} - w_{i,t-1}\rvert,
  \qquad
  \kappa_{i,t} \;=\; \frac{1}{10^4}\left(3\Big(\frac{11}{\log_{10}\Mff}\Big)^{6} + 3\right),
\end{equation}
where $w_{i,t}$ is the name's signed portfolio weight and $\kappa_{i,t}$ is the one-way cost
as a fraction of notional. A position that is opened and later closed thus pays $\kappa$ twice
over its life --- once each way --- but nothing in the months it is simply held. Recomputing
the long/short Sharpe net of this cost isolates which factor effects survive realistic
trading.

\subsection{Experiment 2: Software-industry factors and the search for new signals}

On the software universe we test the established software-industry factors (Section 2.2)
using the same evaluation as Experiment 1 (quintile sorts, cross-sectional regressions, and
gross and cost-adjusted long/short Sharpe), and we use the same machinery to search for
genuinely new factors tailored to software economics.

A key research direction is extrapolating the R\&D activity. The defining input of a
software firm is research and development, yet standard accounting treats it as a period
expense and the standard factor library captures only its \emph{level}. A promising place to 
discover signal is therefore the \emph{behaviour} of R\&D rather
than its size: how fast a firm's research spend is growing, how reliably it converts into
realized profit, how stable that commitment is through cycles, and how it is composed
relative to other discretionary spend. We construct candidate factors along these lines ---
R\&D growth, R\&D-to-profit conversion, the stability of research intensity --- and subject each to the 
Experiment~1 evaluation as well as correlation check with both the general market factors and the
established software factors above.

\subsection{Experiment 3: Multivariate model, alpha, and risk-based sizing}

We normalize returns into industry-relative space by subtracting, for each month, that
month's cross-sectional industry-average return from every stock's return. From Experiments
1 and 2 we select the factors that are statistically significant and survive costs, and fit
a multivariate regression of industry-relative return on these factors, pooled across
stocks and time, to produce a combined return prediction $\hat{r}_{i,t}$.

We test the model out-of-sample with a walk-forward scheme: estimate coefficients on data
through month $t$, predict month $t+1$, and form positions on a cost-adjusted basis using
each name's round-trip cost $2\kappa_{i,t}$:
\begin{itemize}
  \item \textbf{Long leg:} rank stocks by $\hat{r}_{i,t+1} - 2\kappa_{i,t}$ and long the top
  quintile --- the names with the highest net expected return.
  \item \textbf{Short leg:} rank stocks by $-\hat{r}_{i,t+1} - 2\kappa_{i,t}$ and short the
  top quintile --- the names whose expected decline most exceeds the cost of shorting them.
\end{itemize}
Both legs are therefore selected only when the expected edge clears the trading cost. We
then characterize the resulting portfolio in three ways.

\textbf{Alpha.} Regress the portfolio's return on the industry-average return,
$r^{\text{port}}_t = \alpha + \beta\, r^{\text{ind}}_t + \varepsilon_t$. The intercept
$\alpha$ is the return not explained by industry exposure --- the alpha the strategy is
after --- while $\beta$ confirms that residual directional industry exposure is negligible.

\textbf{Internal risk.} Compute the average pairwise correlation among the portfolio's
holdings; a low value indicates the book's risk is genuinely diversified across independent
positions rather than concentrated in one repeated bet.

\textbf{Volatility sizing.} Replace equal weighting within each leg with inverse-volatility
(risk-based) sizing under a fixed volatility target, and compare against equal weighting;
the configuration with the highest and most stable risk-adjusted return is selected as the
most robust portfolio.

The headline metric remains the out-of-sample, post-cost Sharpe ratio, with turnover,
drawdowns, and stability across sub-periods and regimes reported alongside.

%==================================================================
\begin{thebibliography}{99}
\small
\bibitem[Bernard \& Thomas(1989)]{bernard1989} Bernard, V., Thomas, J. (1989). Post-earnings-announcement drift. Journal of Accounting Research 27.
\bibitem[Cooper et al.(2008)]{cooper2008} Cooper, M., Gulen, H., Schill, M. (2008). Asset growth and the cross-section of stock returns. Journal of Finance 63(4).
\bibitem[Fama \& French(1992)]{fama1992} Fama, E., French, K. (1992). The cross-section of expected stock returns. Journal of Finance 47(2).
\bibitem[Fama \& MacBeth(1973)]{famamacbeth1973} Fama, E., MacBeth, J. (1973). Risk, return, and equilibrium: empirical tests. Journal of Political Economy 81(3).
\bibitem[Frazzini \& Pedersen(2014)]{frazzini2014} Frazzini, A., Pedersen, L. (2014). Betting against beta. Journal of Financial Economics 111(1).
\bibitem[Grinold(1989)]{grinold1989} Grinold, R. (1989). The fundamental law of active management. Journal of Portfolio Management 15(3).
\bibitem[Jegadeesh(1990)]{jegadeesh1990} Jegadeesh, N. (1990). Evidence of predictable behavior of security returns. Journal of Finance 45(3).
\bibitem[Jegadeesh \& Titman(1993)]{jegadeesh1993} Jegadeesh, N., Titman, S. (1993). Returns to buying winners and selling losers. Journal of Finance 48(1).
\bibitem[Novy-Marx(2013)]{novymarx2013} Novy-Marx, R. (2013). The other side of value: the gross profitability premium. Journal of Financial Economics 108(1).
\bibitem[Peters \& Taylor(2017)]{petersTaylor2017} Peters, R., Taylor, L. (2017). Intangible capital and the investment-q relation. Journal of Financial Economics 123(2).
\bibitem[Pontiff \& Woodgate(2008)]{pontiff2008} Pontiff, J., Woodgate, A. (2008). Share issuance and cross-sectional returns. Journal of Finance 63(2).
\bibitem[Sloan(1996)]{sloan1996} Sloan, R. (1996). Do stock prices fully reflect information in accruals and cash flows? The Accounting Review 71(3).
\end{thebibliography}

\end{document}
